% AUTO-GENERATED by paper/gen_tables.py -- do not edit by hand; edit the data and re-run.
\begin{table}[t]
\centering
\caption{Static analysis: attack verdicts predicted from the configurations (majority over locations). A star marks a verdict that differs from the design-analysis expectation in Table~\ref{tab:attacks}; the expected value is in parentheses. Computed from the configuration files by a static reachability analyser (\texttt{tools/policy\_analyzer.py}); this is \emph{not} a Packet Tracer measurement.}
\label{tab:static}
\footnotesize
\begin{tabular}{@{}llcccc@{}}
\toprule
ID & Attack & D0 & D1 & D2 & D3 \\
\midrule
A1 & Recon ping sweep & \oS & \oS & \oP & \oP \\
A2 & EHR access (HTTP) & \oS & \oS & \oS & \oB \\
A3 & IoT-server / cross-class & \oS & \oS & \oS & \oP \\
A4 & Telnet/SSH to infra & \oS & \oS & \oB & \oB \\
A5 & Rogue DHCP server & \oS & \oS & \oS & \oB \\
A6 & Unauthorised device & \oS & \oS & \oS & \oB \\
A7 & Guest $\rightarrow$ EHR pivot & \oS & \oS & \oB & \oB \\
\bottomrule
\end{tabular}
\end{table}

\begin{table}[t]
\centering
\caption{Static analysis: hosts and services reachable from the compromised node, mean over locations. ``Any address'' also counts a device reached through any of its interface addresses; ``listed'' counts only the addresses in the test plan. Computed from the configuration files by a static reachability analyser (\texttt{tools/policy\_analyzer.py}); this is \emph{not} a Packet Tracer measurement.}
\label{tab:staticsurface}
\footnotesize
\begin{tabular}{@{}lcccc@{}}
\toprule
& \multicolumn{2}{c}{Hosts} & \multicolumn{2}{c}{Services} \\
\cmidrule(lr){2-3}\cmidrule(l){4-5}
Design & any addr. & listed & any addr. & listed \\
\midrule
D0 Flat & 23 & 23 & 10 & 10 \\
D1 VLAN-only & 23 & 23 & 10 & 10 \\
D2 Zone ACL & 8 & 8 & 4 & 4 \\
D3 Micro-seg & 2.7 & 2.7 & 1 & 1 \\
\bottomrule
\end{tabular}
\end{table}
